\scanpage{61}\mspage{9}
% archive foliation: 60
$x$ pt de $I\wedge_{{}_2E_0}({}_4E_0)\setminus I$

On trouve ainsi
\[
\tag{27}\Delta=
\underbrace{\left(\underbrace{\underbrace{I\wedge_{\underline M_0}\underline M}_{\text{degré }16}\setminus\underbrace{I}_{\text{degré }4}}_{\text{degré }12}\right)/\{1,\sigma\}}_{\text{degré }2}
\]
où on a posé maintenant
\[
\tag{28}\left\{\begin{array}{l}
\underline M\simeq{}_4E_0\quad\text{schéma en groupes loc. isom. à }(\mathbb Z/4\mathbb Z)^2_S\\
\underline M_0\simeq\underline M\otimes_{\mathbb Z/4\mathbb Z}\mathbb Z/2\simeq2\underline M\simeq{}_2E_0=\underline V
\end{array}\right.
\]
Ainsi \inserted{(dans le cas $E=E_0$, disons)} les pts d'ordre $4$ de $E_0$ se groupent par $6$ paires d'éléments symétriques, qui \unsure{fournit|fournissent} la configuration octaédrale $\Delta$ dans $X$.

On trouve ainsi le diagramme (presque cartésien)
\[
\tag{29}
\diagram{\begin{tikzcd}[ampersand replacement=\&,column sep=large,row sep=huge]
\underline I\arrow[r,phantom,"\subset"]\&E\arrow[r,"\text{étale}"]\arrow[d,"\substack{\text{ramifié en}\\\text{\unsure{dessus} des pts}\\\text{de }\underline I}"]\&E_0\;(\simeq E/\underline M_0)\arrow[d,"\substack{\text{ramifié en dessus}\\\text{des pts de }D_J\cup t(S)}"]\\
\underline I\arrow[r,hook]\&X\arrow[r,"\substack{\text{ramifié en}\\\text{dessus des}\\\text{pts de }D_J}"]\&\Sigma_J\\
\&\Delta\arrow[u,phantom,"\rotatebox{90}{$\subset$}"]\arrow[r,dashed]\&D_J\arrow[u,phantom,"\rotatebox{90}{$\subset$}"]
\end{tikzcd}}
\]
qui est cartésien au dessus de $\bar\Sigma_J^*=\bar\Sigma_J\setminus D_J$
